\documentclass[12pt]{article}
%\include{amsfonts}
\usepackage{amssymb,amsmath,amsthm,latexsym,epsfig,euscript,multicol}
\usepackage{enumitem}
\usepackage[utf8x]{inputenc}
\usepackage{booktabs}
\usepackage{babel}[spanish]
% \usepackage{enumerate}

% Caracteres especiales
\def\A{\mathbb{A}}
\def\C{\mathbb{C}}
\def \N{\mathbb{N}}
\def \P{\mathbb{P}}
\def \Q{\mathbb{Q}}
\def \R{\mathbb{R}}
\def \Z{\mathbb{Z}}
\def \sen{\textrm{sen}}

\theoremstyle{definition}
\newtheorem{ejer}{Ejercicio}
\newcommand{\bej}{\begin{ejer}}
\newcommand{\fej}{\end{ejer}}

\def\dt{\Delta t}
\def\dx{\Delta x}

\topmargin-2cm \vsize 29.5cm \hsize 21cm
\setlength{\textwidth}{16.75cm}\setlength{\textheight}{23.5cm}
\setlength{\oddsidemargin}{0.0cm}
\setlength{\evensidemargin}{0.0cm}


\begin{document}

%----------------------------------------------------------
%Encabezado:

\centerline{{\small Universidad de Buenos Aires - Facultad de Ciencias Exactas y Naturales - Depto. de Matemática}}

\vskip 0.2cm

\hrule

\vskip 0.2cm

 \centerline{{\bf\Large{\sc Investigación Operativa}}}

 \vskip 0.2cm

 \centerline{\ttfamily Segundo Cuatrimestre 2025}

\vskip 0.2cm

 \hrule

 \bigskip
 \centerline{\bf Práctica 3: Dualidad y Branch \& Bound}
 \bigskip

%--------------------------------------------------------- 
% 1
\bej
    Formular el dual de los siguientes problemas lineales:
    \begin{multicols}{2}
    \begin{enumerate}[label=\alph*)]
        \item \[
            \begin{array}{rrcl}
            \min & \multicolumn{2}{l}{z = x_1+4x_2+x_3} \\
            s.a: & 2x_1-2x_2+x_3 & = & 4 \\\
             & x_1-x_3 & = & 1 \\
             & x_2, x_3& \geq & 0 \\
             & x_1 & \multicolumn{2}{l}{\text{libre}}
            \end{array}
            \]
        \item \[
                \begin{array}{rrcl}
                \max & \multicolumn{2}{l}{z = -3x_1-9x_2+x_3-x_5} \\
                s.a: & -5x_1+2x_2-x_4 & \leq & 12 \\
                 & -3x_1+x_2+2x_3+x_5 & \leq & 7 \\
                 & x_1 - 3x_4 + x_5 & \geq & -4 \\
                  & x_1,x_2,x_3 & \geq & 0 \\
                  & x_4, x_5 & \leq & 0
                \end{array}
                \]
    \end{enumerate}
    \end{multicols}
\fej

%---------------------------------------------------------
\bej
    Probar que el problema dual del problema dual es el problema primal.
\fej

%---------------------------------------------------------
\bej
    Encuentre un problema no factible cuyo dual tampoco sea factible.
\fej

%---------------------------------------------------------
\bej
    Dado el siguiente problema:
    \[
    \begin{array}{rrcl}
    \min & \multicolumn{3}{l}{z = x_1+x_2+3x_3+2x_4+8x_5+2x_6+x_7} \\
    s.a: & 2x_1+2x_2+x_3+32x_4+3x_5+20x_6+11x_7 & \geq & 1 \\
     & x & \geq & 0
    \end{array}
    \]
    \begin{enumerate}[label=\alph*)]
        \item Hallar el valor óptimo de la función objetivo.
        \item Hallar el punto donde se alcanza a partir de la resolución del problema dual.
    \end{enumerate}
\fej

%---------------------------------------------------------
\bej 
    Resolver cada uno de los siguientes problemas a partir del planteo y resolución de su correspondiente problema dual:
    \begin{multicols}{2}
    \begin{enumerate}[label=\alph*)]
        \item \[
            \begin{array}{rrcl}
            \min & \multicolumn{2}{l}{z = -5x_1-7x_2-12x_3+x_4} \\
            s.a: & 2x_1+3x_2+2x_3+x_4 & \leq & 38 \\
             & 3x_1+2x_2+4x_3-x_4 & \leq & 55 \\
             & x& \geq & 0
            \end{array}
            \]
    \item \[
            \begin{array}{rrcl}
            \max & \multicolumn{2}{l}{z = 5x_1+3x_2+2x_3} \\
            s.a: & 4x_1+5x_2+2x_3+x_4 & \leq & 20 \\
             & 3x_1+4x_2-x_3+x_4 & \leq & 30 \\
             & x& \geq & 0
            \end{array}
            \]
    \end{enumerate}
    \end{multicols} 
\fej


\bej Para cada uno de los siguientes modelos, decidir si $x^\ast$ es una solución óptima utilizando el problema dual.
    \begin{enumerate}
        \item 
        \[x^\ast = (0, 0, 8, 3) \]
        \[
        \begin{array}{rrrrrrrrcl}
        \max & 3x_1&+&2x_2&-&x_3&+&3x_4 & &  \\
        s.a: & x_1&+&2x_2&-&x_3&+&2x_4 & \leq & -2 \\
         & 3x_1&-&x_2&+&x_3&-&x_4 & \leq & 5 \\
         & -2x_1&+&x_2&-&3x_3&+&3x_4 & \leq & 10 \\
         & \multicolumn{7}{r}{x_1,x_2,x_3,x_4} & \geq & 0
        \end{array}
        \]
        \item     
        \[x^\ast = (0, 2, 0, 7, 0)\]
        \[\begin{array}{rrrrrrrrrrrr}
        \max & 8x_1 & - & 9x_2 & + & 12x_3 & + & 4x_4 & + & 11x_5   \\
        \text{s.a:} & 2x_1 & - & 3x_2 & + &4x_3 & + & x_4 & + & 3x_5 & \leq & 1 \\
         & x_1 & + & 7x_2 & + &3x_3 & - & 2x_4 & + & x_5 & \leq & 1 \\
         & 5x_1 & + & 4x_2 & - & 6x_3 & + & 2x_4 & + & 3x_5 & \leq & 22 \\
         & \multicolumn{9}{r}{x_1,x_2,x_3,x_4,x_5} & \geq & 0
    \end{array}
    \]

    \end{enumerate}
\fej


%---------------------------------------------------------
%\bej
%    Considere el siguiente problema de programación lineal:
%    \[
%    \begin{array}{rrcl}
%    \max & \multicolumn{3}{l}{z = c_1x_1+c_2x_2} \\
%    s.a: & 2x_1+3x_2 & \leq & b_1 \\
%     & x_1-x_2 & \leq & b_2 \\
%     & x_1,x_2 & \geq & 0
%    \end{array}
%    \]
%    con $c_1=1,\;c_2=3,\;b_1=6$ y $b_2=2$.
%    \begin{enumerate}[label=\alph*)]
%        \item Aplique el algoritmo Simplex para encontrar el óptimo.
%        \item ¿En qué intervalo puede moverse $c_1$ y que el óptimo siga siendo el mismo?
%        \item ¿En qué intervalo puede moverse $b_1$ y que el óptimo siga siendo el mismo?
%        \item Formule el problema dual del problema original y encuentre el óptimo del dual utilizando el Teorema de Holgura Complementaria.
%        \item Suponga que $b_1$ aumenta de $6$ a $6.5$, ¿en cuánto mejora la función objetivo?
%        \item Suponga que $b_2$ aumenta de $2$ a $3$, ¿en cuánto mejora la función objetivo?
%    \end{enumerate}
%\fej
%
%%---------------------------------------------------------
%\bej
%    Sea $\theta\in\mathbb{R}$, considere el problema:
%    \[
%    \begin{array}{rrcl}
%    \max & \multicolumn{3}{l}{z = (10-4\theta)x_1+(4-\theta)x_2+(7+\theta)x_3} \\
%    s.a: & 3x_1+x_2+2x_3 & \leq & 7 \\
%     & 2x_1+x_2+3x_3 & \leq & 5 \\
%     & x & \geq & 0
%    \end{array}
%    \]
%\fej
%
%\begin{enumerate}[label=\alph*)]
%    \item Aplique el algoritmo Simplex para encontrar el óptimo con $\theta = 0$.
%    \item Determine el intervalo de valores de $\theta$ para que la solución básica factible obtenida en el ítem anterior siga siendo óptima.
%    \item Sea $\theta$ perteneciente al intervalo hallado en el ítem anterior, encuentre el intervalo de valores de $b_1$ tales que la base factible óptima siga siendo la misma. Repetir el procedimiento para $b_2$.
%    \item Analice el caso en el que $b_1$ y $b_2$ se mueven a la vez.
%    \item Formule el dual del problema original.
%    \item Para $\theta=0$, hallar el óptimo del problema dual utilizando el Teorema de Holgura Complementaria.
%\end{enumerate}

%---------------------------------------------------------
\bej
    Resuelva utilizando Branch \& Bound los siguientes problemas:
    \begin{multicols}{2}
    \begin{enumerate}[label=\alph*)]
    \item \[
        \begin{array}{rrcl}
        \max & \multicolumn{3}{l}{z = 2x_1+x_2} \\
        s.a: & 2x_1+3x_2 & \leq & 19 \\
         & 7x_1+3x_2 & \leq & 43 \\
         & x_1, x_2 & \geq & 0 \\
         & x_1, x_2 &\in & \mathbb{Z}
        \end{array}
        \]

    \item \[
        \begin{array}{rrcl}
        \min & \multicolumn{3}{l}{z = -10x_1-15x_2} \\
        s.a: & 8x_1+4x_2 & \leq & 40 \\
         & 15x_1+30x_2 & \leq & 20 \\
          & x_1, x_2 & \geq & 0 \\
         & x_1, x_2 &\in & \mathbb{Z}
        \end{array}
        \]
    \item \[
        \begin{array}{rrcl}
        \min & \multicolumn{3}{l}{z = -4x_1+5x_2} \\
        s.a: & 5x_1-2x_2 & \geq & -2 \\
         & 2x_1-4x_2 & \leq & 5 \\
         & x_1 & \leq & 5 \\
         & x_2 & \leq & 4 \\
           & x_1, x_2 & \geq & 0 \\
         & x_1, x_2 &\in & \mathbb{Z}

        \end{array}
        \]
    \end{enumerate}
    \end{multicols} 
\fej

\bej
    Para los ejemplos del ejercicio anterior, supongamos que por restricciones de tiempo, no
    podemos expandir el árbol más allá del segundo nivel. Elija la mejor solución al problema en
    ese caso y dé una medida de la calidad de la solución encontrada.
\fej
\end{document}
%--------------------------------------------------------- 